\cleardoublepage
\phantomsection
\begin{planbox}[แผนบริหารการสอนประจำบทที่ 3\quad กฎการเคลื่อนที่ของนิวตันและแรงหนืด (สัปดาห์ที่ 3)]
\noindent\textbf{1. หัวข้อเนื้อหาประจำบท (บรรยาย 2 ชั่วโมง และปฏิบัติการ 3 ชั่วโมง)}
\begin{itemize}
    \item 3.1 กฎการเคลื่อนที่ 3 ข้อของนิวตันในมุมมองชีวฟิสิกส์ระดับเซลล์
    \item 3.2 แรงต้านการเคลื่อนที่ในของเหลว ความหนืด และกฎของสโตกส์
    \item 3.3 อัตราเร็วสุดท้ายของการตกตะกอนของเซลล์จุลินทรีย์ในอาหารเลี้ยงเชื้อเหลว
    \item 3.4 ปฏิบัติการที่ 3 กฎของนิวตันและการวัดความหนืดของสารละลายชีวภาพด้วยวิธีปล่อยทรงกลมตกในของเหลว
\end{itemize}

\noindent\textbf{2. วัตถุประสงค์เชิงพฤติกรรมและผลลัพธ์การเรียนรู้รายวิชา (CLOs ที่เชื่อมโยง PLOs)}
\begin{enumerate}
    \item อธิบายกฎการเคลื่อนที่ 3 ข้อของนิวตันและใช้สมการ $F = ma$ คำนวณความสัมพันธ์ระหว่างแรง มวล และความเร่งได้ \textbf{(CLO1 เชื่อมโยง PLO2)}
    \item อธิบายอิทธิพลของความหนืด ($\eta$) และขนาดรัศมีเซลล์ ($r$) ต่ออัตราเร็วการตกตะกอนตามกฎของสโตกส์ได้ \textbf{(CLO1, CLO2 เชื่อมโยง PLO2, PLO3)}
    \item ปฏิบัติการทดลองวัดความหนืดของของเหลวอย่างระมัดระวังและสื่อสารผลการทดลองเป็นกราฟและรายงานได้อย่างถูกต้อง \textbf{(CLO3, CLO4 เชื่อมโยง PLO1, PLO5, PLO6)}
\end{enumerate}

\noindent\textbf{3. กิจกรรมการจัดการเรียนรู้เชิงรุก (Active Learning \& OBE)}
\begin{itemize}
    \item \textbf{Demonstration \& Inquiry}\quad เปรียบเทียบการตกของเม็ดบีดส์ในน้ำเปล่ากับในกลีเซอรีน (ซึ่งใช้ผสมเก็บรักษาเชื้อจุลินทรีย์ที่ $-80\,^\circ\text{C}$) เพื่อให้นักศึกษาเห็นภาพแรงหนืดอย่างชัดเจน
    \item \textbf{ฝึกวิเคราะห์ความสัมพันธ์เชิงสัดส่วน}\quad ชี้ให้เห็นว่าอัตราเร็วตกตะกอนแปรผันตรงกับกำลังสองของรัศมีเซลล์ ($v_t \propto r^2$) จึงอธิบายได้ว่าทำไมยีสต์จึงตกตะกอนเร็วกว่าแบคทีเรียมาก
\end{itemize}

\noindent\textbf{4. สื่อการเรียนรู้และแหล่งเรียนรู้}
\begin{itemize}
    \item เอกสารคำสอนบทที่ 3 แผนภาพสมดุลแรง 3 แรงบนเซลล์จุลินทรีย์ และชุดกระบอกตวงวัดความหนืดของเหลว
\end{itemize}

\noindent\textbf{5. การวัดและประเมินผลประจำบท}
\begin{itemize}
    \item คำถามท้ายคาบเรื่องสมดุลแรงและกฎของสโตกส์ \textbf{(CLO1, CLO2)}
    \item การประเมินทักษะการทดลองและรายงานปฏิบัติการที่ 3 \textbf{(CLO3, CLO4)}
\end{itemize}
\end{planbox}
\clearpage

\chapter{กฎการเคลื่อนที่ของนิวตันและแรงหนืด}
\index{กฎการเคลื่อนที่ของนิวตัน}\index{ความหนืด}\index{กฎของสโตกส์}\index{อัตราเร็วสุดท้าย}

\section{กฎการเคลื่อนที่สามข้อของนิวตันในมุมมองชีวฟิสิกส์}

เซอร์ไอแซก นิวตัน ได้สรุปกฎพื้นฐานที่ควบคุมการเคลื่อนที่ของวัตถุทุกขนาดตั้งแต่ระดับดวงดาวไปจนถึงระดับเซลล์จุลินทรีย์ไว้ 3 ข้อ ซึ่งสามารถทำความเข้าใจโดยไม่ต้องใช้คณิตศาสตร์ซับซ้อน ดังนี้

\begin{enumerate}
    \item \textbf{กฎข้อที่ 1 กฎของความเฉื่อย} (\textbf{Law of Inertia})\quad เมื่อผลรวมของแรงลัพธ์ที่กระทำต่อวัตถุมีค่าเป็นศูนย์ ($\sum F = 0$) วัตถุที่อยู่นิ่งก็จะอยู่นิ่งต่อไป ส่วนวัตถุที่กำลังเคลื่อนที่ก็จะเคลื่อนที่ด้วยความเร็วคงที่ ในระดับจุลินทรีย์ เนื่องจากเซลล์แบคทีเรียมีมวลน้อยมาก (ประมาณ $10^{-15}\text{ kg}$) แรงต้านจากความหนืดของน้ำจึงมีอิทธิพลเหนือความเฉื่อยอย่างสิ้นเชิง ทันทีที่แบคทีเรียหยุดหมุนแฟลกเจลลา มันจะหยุดเคลื่อนที่แทบจะในทันทีภายในระยะทางไม่ถึงขนาดอะตอม
    \item \textbf{กฎข้อที่ 2 กฎของความเร่ง} (\textbf{Law of Acceleration})\quad เมื่อมีแรงลัพธ์ที่ไม่เป็นศูนย์กระทำต่อวัตถุมวล $m$ วัตถุจะเกิดความเร่ง $a$ ในทิศทางเดียวกับแรงลัพธ์ ตามสมการที่ (3.1)
    \begin{equation}
        F_{\text{net}} = m a
    \end{equation}
    โดยแรงมีหน่วยเป็น\textbf{นิวตัน} ($\text{N}$), มวลมีหน่วยเป็นกิโลกรัม ($\text{kg}$) และความเร่งมีหน่วยเป็นเมตรต่อวินาทียกกำลังสอง ($\text{m/s}^2$)
    \item \textbf{กฎข้อที่ 3 กฎของแรงกิริยาและแรงปฏิกิริยา} (\textbf{Action equals Reaction})\quad ทุกแรงกิริยาจะมีแรงปฏิกิริยาที่มีขนาดเท่ากันแต่มีทิศทางตรงกันข้ามเสมอ ($F_{\text{action}} = -F_{\text{reaction}}$) ตัวอย่างเช่น เมื่อแฟลกเจลลาของแบคทีเรียหมุนผลักน้ำไปด้านหลัง (แรงกิริยา) น้ำจะออกแรงปฏิกิริยาผลักตัวเซลล์แบคทีเรียให้พุ่งไปข้างหน้า
\end{enumerate}

\section{แรงหนืด กฎของสโตกส์ และการตกตะกอนของเซลล์}

เมื่อเซลล์จุลินทรีย์ทรงกลมรัศมี $r$ เคลื่อนที่ด้วยอัตราเร็ว $v$ ผ่านอาหารเลี้ยงเชื้อเหลวที่มี\textbf{สัมประสิทธิ์ความหนืด} (\textbf{Viscosity}, $\eta$ อ่านว่า อีตา มีหน่วยเป็น $\text{Pa}\cdot\text{s}$) ของเหลวจะออกแรงต้านการเคลื่อนที่ซึ่งเรียกว่า \textbf{แรงหนืดตามกฎของสโตกส์} (\textbf{Stokes' Drag Force}) ดังสมการที่ (3.2)

\begin{equation}
    F_{\text{drag}} = 6 \pi \eta r v
\end{equation}

เมื่อเราตั้งหลอดอาหารเลี้ยงเชื้อเหลวทิ้งไว้นิ่งๆ เซลล์จุลินทรีย์จะค่อยๆ จมลงสู่ก้นหลอดภายใต้แรงกระทำ 3 แรง ได้แก่ \textbf{แรงโน้มถ่วง} ($F_g$) ดึงลงด้านล่าง, \textbf{แรงลอยตัว} ($F_b$) ดันขึ้นด้านบน และ\textbf{แรงหนืด} ($F_{\text{drag}}$) ต้านการจมดันขึ้นด้านบน ดังแสดงใน\mbox{ภาพที่} 3.1 เมื่อแรงขึ้นรวมเท่ากับแรงลงพอดี ($\sum F = 0$) เซลล์จะจมลงด้วย\textbf{อัตราเร็วสุดท้ายคงที่} (\textbf{Terminal Velocity} หรืออัตราเร็วตกตะกอน $v_t$) ตามสมการที่ (3.3)

\begin{equation}
    v_t = \frac{2 r^2 (\rho_p - \rho_f) g}{9 \eta}
\end{equation}

โดยที่ $\rho_p$ คือความหนาแน่นของเซลล์ ($\text{kg/m}^3$), $\rho_f$ คือความหนาแน่นของของเหลว ($\text{kg/m}^3$) และ $g \approx 9.8\text{ m/s}^2$ คือความเร่งโน้มถ่วงของโลก

\begin{figure}[H]
\centering
\resizebox{0.84\textwidth}{!}{%
\begin{tikzpicture}[font=\sffamily, >=Stealth]
    % Liquid Cylinder
    \fill[boxBlueBg] (2.0,0.4) rectangle (8.0,5.2);
    \draw[physRoyalBlue, line width=1.6pt] (2.0,5.4) -- (2.0,0.4) -- (8.0,0.4) -- (8.0,5.4);
    \draw[physCyan, line width=1.2pt, decorate, decoration={snake, amplitude=1.2mm, segment length=8mm}] (2.0,5.2) -- (8.0,5.2);

    % Spherical Microbial Cell
    \shade[ball color=physAmber!85] (5.0,2.8) circle (0.75cm);
    \node[white, font=\bfseries] at (5.0,2.8) {เซลล์ $r, \rho_p$};

    % Forces
    \draw[->, line width=2.0pt, color=crimsonAlert] (5.0,2.05) -- (5.0,0.75) node[right, font=\small\bfseries] {แรงโน้มถ่วง $F_g$ (ดึงลง)};
    \draw[->, line width=2.0pt, color=bioEmerald] (4.65,3.55) -- (4.65,4.75) node[left, font=\small\bfseries] {แรงหนืด $F_{\text{drag}} = 6\pi\eta r v_t$};
    \draw[->, line width=2.0pt, color=physCyan] (5.35,3.55) -- (5.35,4.45) node[right, font=\small\bfseries] {แรงลอยตัว $F_b$};

    % Explanation Box on Right
    \node[anchor=west, align=left, draw=physRoyalBlue, fill=softSmoke, rounded corners=3pt, inner sep=8pt, text width=6.2cm] at (8.8,2.8) {
        \textbf{\color{physRoyalBlue}สมดุลแรงขณะตกตะกอน ($\sum F = 0$)}\\[4pt]
        {\small แรงขึ้นรวม $=$ แรงลง}\\
        {\small $F_{\text{drag}} + F_b = F_g$}\\[4pt]
        {\footnotesize\color{charcoal} สังเกตว่า $v_t \propto r^2$ ดังนั้นหากเซลล์มีรัศมีใหญ่ขึ้น $5$ เท่า อัตราเร็วการตกตะกอนจะเพิ่มขึ้นถึง $5^2 = 25$ เท่า!}
    };
\end{tikzpicture}%
}
\caption{แผนภาพวัตถุอิสระแสดงสมดุลของแรงทั้งสามที่กระทำต่อเซลล์จุลินทรีย์ขณะจมลงด้วยอัตราเร็วสุดท้ายในของเหลว}
\end{figure}

\begin{bioappbox}[ทำไมเซลล์ยีสต์จึงตกตะกอนนอนก้นหลอดเร็วกว่าเซลล์แบคทีเรียถึง 25 เท่า]
จากสมการที่ (3.3) อัตราเร็วตกตะกอนแปรผันตรงกับกำลังสองของรัศมีเซลล์ ($v_t \propto r^2$) เซลล์ยีสต์มีรัศมีเฉลี่ยประมาณ $r_{\text{yeast}} \approx 2.5\text{ }\mu\text{m}$ ในขณะที่เซลล์แบคทีเรียทรงกลมมีรัศมีประมาณ $r_{\text{bacteria}} \approx 0.5\text{ }\mu\text{m}$ (ต่างกัน $5$ เท่า) เมื่อยกกำลังสองจะได้ $(5)^2 = 25$ เท่า ทำให้เซลล์ยีสต์ตกตะกอนลงก้นหลอดทดลองได้เองภายในไม่กี่ชั่วโมง ในขณะที่แบคทีเรียต้องอาศัยเครื่องปั่นเหวี่ยงความเร็วสูงเพื่อช่วยเร่งการตกตะกอน
\end{bioappbox}

\begin{worked}[การเปรียบเทียบอัตราเร็วตกตะกอนและการคำนวณแรงลัพธ์]
\textbf{โจทย์ปัญหา}\quad สปอร์ของเชื้อราชนิดหนึ่งมีรัศมี $r_A = 3\text{ }\mu\text{m}$ และสปอร์ของแบคทีเรียมีรัศมี $r_B = 1\text{ }\mu\text{m}$ หากสปอร์ทั้งสองชนิดมีความหนาแน่นเท่ากันและแขวนลอยอยู่ในอาหารเหลวชนิดเดียวกัน สปอร์ของเชื้อราจะจมลงก้นหลอดเร็วกว่าสปอร์ของแบคทีเรียกี่เท่า

\textbf{วิธีทำทีละขั้นตอนฉบับเข้าใจง่าย}
\begin{enumerate}
    \item \textbf{ขั้นที่ 1 ใช้หลักสัดส่วนตรงจากสมการกฎของสโตกส์ ($v_t \propto r^2$)}
    \[
        \frac{v_A}{v_B} = \left(\frac{r_A}{r_B}\right)^2
    \]
    \item \textbf{ขั้นที่ 2 แทนค่ารัศมีแล้วคำนวณกำลังสอง}
    \[
        \frac{v_A}{v_B} = \left(\frac{3\text{ }\mu\text{m}}{1\text{ }\mu\text{m}}\right)^2 = 3^2 = 9\text{ เท่า}
    \]
\end{enumerate}
\textbf{คำตอบ}\quad สปอร์ของเชื้อราจะตกตะกอนเร็วกว่าสปอร์ของแบคทีเรียเท่ากับ $9$ เท่า
\end{worked}

\section{ปฏิบัติการประจำสัปดาห์ที่ 3}

\begin{labbox}[ปฏิบัติการที่ 3\quad กฎของนิวตันและการวัดความหนืดของของเหลวชีวภาพ (3 ชั่วโมง)]
\noindent\textbf{ผลลัพธ์การเรียนรู้ที่ประเมิน}\quad CLO1 (กฎของนิวตันและแรงหนืด), CLO2 (เชื่อมโยงความหนืดกับการตกตะกอนในงานจุลชีววิทยา), CLO3 (ความปลอดภัยและวินัยในห้องแล็บ)

\vspace{4pt}
\noindent\textbf{1. วัสดุและอุปกรณ์}
\begin{itemize}
    \item กระบอกตวงทรงสูงขนาด $500\text{ mL}$ จำนวน 2 ใบ บรรจุน้ำกลั่นและสารละลายกลีเซอรีน (หรือน้ำเชื่อมเข้มข้นจำลองอาหารเลี้ยงเชื้อความหนืดสูง)
    \item ลูกเหล็กทรงกลมขนาดเล็ก (ทราบรัศมี $r$ และความหนาแน่น $\rho_p$), ไมโครมิเตอร์, แม่เหล็กสำหรับดึงลูกเหล็ก และนาฬิกาจับเวลา
\end{itemize}

\noindent\textbf{2. ขั้นตอนการทดลอง}
\begin{enumerate}
    \item ติดแถบยางรัดรอบกระบอกตวงที่ขีดบอกระดับบน (ต่ำกว่าผิวของเหลว $5\text{ cm}$ เพื่อให้ลูกเหล็กเข้าสู่ความเร็วคงที่ก่อน) และขีดบอกระดับล่าง ห่างกันเป็นระยะทาง $s = 20\text{ cm}$ ($0.20\text{ m}$)
    \item ปล่อยลูกเหล็กทรงกลมลงตรงกลางกระบอกตวง จับเวลา $t$ ที่ลูกเหล็กเคลื่อนที่ผ่านระหว่างแถบยางทั้งสองเส้น ทำซ้ำ 3 ครั้ง คำนวณอัตราเร็วสุดท้าย $v_t = s/t$ และคำนวณสัมประสิทธิ์ความหนืด $\eta$ จากสมการที่ (3.3)
\end{enumerate}

\noindent\textbf{3. ตารางบันทึกผลการทดลอง}
\begin{center}
\footnotesize\setlength{\tabcolsep}{3.5pt}
\begin{tabularx}{\linewidth}{>{\raggedright\arraybackslash}X c c c c c}
\toprule
\textbf{ชนิดของเหลว} & \textbf{ครั้งที่ 1 ($\text{s}$)} & \textbf{ครั้งที่ 2 ($\text{s}$)} & \textbf{เฉลี่ย $t$ ($\text{s}$)} & \textbf{$v_t$ ($\text{m/s}$)} & \textbf{ความหนืด $\eta$} \\
\midrule
น้ำกลั่น (ควบคุม) & ............ & ............ & ............ & ............ & ............ \\
กลีเซอรีนเข้มข้น & ............ & ............ & ............ & ............ & ............ \\
\bottomrule
\end{tabularx}
\end{center}
\end{labbox}

\section*{คำถามท้ายบทที่ 3 (ทบทวนแนวคิดและการคำนวณเบื้องต้น)}
\begin{enumerate}
    \item หากเราเพิ่มอุณหภูมิของอาหารเลี้ยงเชื้อเหลวให้สูงขึ้น ความหนืด ($\eta$) ของของเหลวจะลดลง จงอธิบายว่าอัตราเร็วการตกตะกอน ($v_t$) ของเซลล์จุลินทรีย์จะเพิ่มขึ้นหรือลดลง เพราะเหตุใด
    \item เซลล์สาหร่ายเซลล์เดียวชนิดหนึ่งมีรัศมี $6\text{ }\mu\text{m}$ ในขณะที่เซลล์แบคทีเรียมีรัศมี $1.5\text{ }\mu\text{m}$ จงหาว่าเซลล์สาหร่ายจะตกตะกอนเร็วกว่าเซลล์แบคทีเรียกี่เท่า
\end{enumerate}
